\subsection{The main theorems}
\label{series:III:main-theorems}

We fix the marked $(3,4,\infty)$ family of Paper~I, its positive
normal-circle orientation, ordered integral lattice, finite affine twists,
and prescribed cusp fan. The first four theorems use the standard
sphere twist data $|p|=1$. The map $e:X\to S^3$ is the repaired quotient
with its target and collar data in \SeriesRef{mapsI:src:A.2.1};
$f:X\to\PP^1$ is its holomorphic algebraic reduction. Every further
elliptic profile carries the source section, vanishing linearization,
lattice and affine cap data specified in its statement. The explicit
transport dictionary is \ref{sec:III:conventions}.

\hypertarget{main-III-A}{}
\paragraph{Theorem A (The prescribed quotient and relative extensions).}
\label{series:III:A}
The quotient of the regular family extends across the cusp through
the explicit target solid torus constructed in Paper~I, with the
target and collar conditions fixing the extension used here.
Agreement only on the exterior permits all twelve elements of
$\pi_6(S^3)$, because the relative pinch of an interior six-ball
adds any chosen element (\SeriesRef{mapsI:src:A.2.1}, \cref{maps:src:A.2.2}).

\textit{Mechanism.} The target solid torus fixes the relative extension, and an interior-ball pinch realizes every unrestricted exterior-agreement class.
\textit{Relation to literature.} The relative pinch and $\pi_6(S^3)=\Z/12$ are known inputs; the explicit receiving form and its counterexample are proved in \cref{maps:src:A.2.2}.
\hyperlink{outline-III-A}{Proof in outline}; full statements: \SeriesRef{maps:src:A.2.2}.

\hypertarget{main-III-B}{}
\paragraph{Theorem B (The quotient and algebraic-reduction homotopy classes).}
\label{series:III:B}
The marked quotient has
$[e]=6\nu'=\eta_3\eta_4\eta_5$, the element of order two in
$\pi_6(S^3)=\Z/12$.  A regular inverse image is a three-torus;
the monodromy calculation identifies its normal framing with its Lie
framing, whose real $e$-invariant is $1/2$
(\cref{maps:src:B.1.1}, \cref{maps:src:B.1.2}).

The holomorphic algebraic reduction $f:X\to\PP^1$ is
null-homotopic.  Its explicit factorization through the circle quotient
contains the normalized power map of degree $3\cdot4=12$ on $S^3$,
which kills $\pi_6(S^3)$ (\cref{maps:src:B.2.1}).

\textit{Mechanism.} The actual normal-frame comparison identifies the Lie-framed regular fibre, and the exact degree-twelve factorization kills its unstable class.
\textit{Relation to literature.} The real Adams invariant and suspension inputs are known: \cite[Theorem~4.14, p.~423]{mapsAPS2} and \cite[Theorem~4.1, Corollaries~4.2--4.3]{mapsToda52}. The full statements identify their applications to the retained maps.
\hyperlink{outline-III-B}{Proof in outline}; full statements: \SeriesRef{maps:src:B.1.1}, \SeriesRef{maps:src:B.1.2}, \SeriesRef{maps:src:B.2.1}.

\hypertarget{main-III-C}{}
\paragraph{Theorem C (The retained-angle cusp map and finite cell model).}
\label{series:III:C}
The cusp collapse retaining the angular coordinate represents
$\eta_4\eta_5\ne0$ in $\pi_6(S^4)$.  At a triple point its
regular inverse image is a two-torus, and the explicit normal-frame
comparison makes this the Lie-framed torus of Arf invariant one
(\cref{maps:src:B.3.1}, \cref{maps:src:B.3.2}, \cref{maps:src:B.3.3}).

The regular torus mapping cylinders, the finite affine caps
and the marked honeycomb cusp fit into a finite CW model with cell
counts $(5,20,45,56,42,17,3)$.  The integral attaching maps give the
homotopy type of the six-sphere
(\cref{maps:src:B.5.1}, \cref{maps:src:B.5.2}).

\textit{Mechanism.} Explicit chart-frame homotopies preserve the Lie framing, and the same marked collar maps give integral cell attachments.
\textit{Relation to literature.} The framed-torus invariant is a known topological input, cited at its exact locator in \cref{maps:src:B.3.1}. The product-collar input is \cite[Theorem~6.1, p.~77]{mapsWaldhausen68}; the marked constructions are the receiving applications.
\hyperlink{outline-III-C}{Proof in outline}; full statements: \SeriesRef{maps:src:B.3.1}, \SeriesRef{maps:src:B.5.1}, \SeriesRef{maps:src:B.5.2}.

\hypertarget{main-III-D}{}
\paragraph{Theorem D (The relative Kummer cover and its global completion).}
\label{series:III:D}
The prescribed degree-$N^2$ relative Kummer cover of the cusp,
for $N>1$, has no extension to a smooth finite cover of the whole
threefold, because its marked monodromy continuation has primitive
ray indices that violate the full-fan smoothness equations.
The explicit semidirect cover and its normal monomial completion
give the finite flat normal completion
(\cref{maps:src:C.1.1}, \cref{maps:src:C.1.2}, \cref{maps:src:C.1.3}, \cref{maps:src:C.1.4}, \cref{maps:src:C.1.5}).

\textit{Mechanism.} Adjacent unimodular fan cones force scalar primitive ray indices, while the retained finite monodromy carries the prescribed lattice outside that scalar system.
\textit{Relation to literature.} The covering-space and normal-monomial inputs are credited at the exact locators in \cref{maps:src:C.1.1,maps:src:C.1.3,maps:src:C.1.5}; the statement retains the prescribed entire cusp component and its actual lattice maps.
\hyperlink{outline-III-D}{Proof in outline}; full statements: \SeriesRef{maps:src:C.1.1}, \SeriesRef{maps:src:C.1.2}, \SeriesRef{maps:src:C.1.3}, \SeriesRef{maps:src:C.1.4}, \SeriesRef{maps:src:C.1.5}.

\hypertarget{main-III-E}{}
\paragraph{Theorem E (The cyclic universal covers).}
\label{series:III:E}
For the original marked family with fundamental group
$\Z/d$, where $d$ is prime to six, the universal cover has
$H_2=\Z^{d-1}\oplus\Z/d$, $H_3=\Z/d$, and Euler number
$2d$.  Its cusp divisor products and finite-coefficient cubic
evaluations determine the complete Zhubr invariant, hence one
oriented diffeomorphism type $M_d$ for each $d$
(\cref{maps:src:D.1.1}, \cref{maps:src:D.1.4}, \cref{maps:src:D.1.6}, \cref{maps:src:D.1.9}, \cref{maps:src:D.1.11}).

\textit{Mechanism.} The cyclic cusp matrix and its finite-coefficient cubic evaluations specify the complete published smooth invariant of the constructed cover.
\textit{Relation to literature.} The starting quotient family is credited to \cite[Theorems~7.17 and 7.22, pp.~57 and 60]{AC26}. The body cites Zhubr\textquotesingle s classification at its exact locator and computes the full receiving invariant, including its torsion data.
\hyperlink{outline-III-E}{Proof in outline}; full statements: \SeriesRef{maps:src:D.1.1}, \SeriesRef{maps:src:D.1.4}, \SeriesRef{maps:src:D.1.6}, \SeriesRef{maps:src:D.1.9}, \SeriesRef{maps:src:D.1.11}.

\hypertarget{main-III-F}{}
\paragraph{Theorem F (The marked elliptic profiles and their integral outputs).}
\label{series:III:F}
The primitive specialization lattices and actual normal-circle-torus
maps classify the marked natural-line sphere outputs on all
twenty-one rank-one Mordell--Weil profiles: $59$ patterns on $50$
initial rows, with twelve profiles of unbounded sphere multiplicity,
four of logarithmic order one and five with no sphere output.
Their integral cokernels and full characteristic evaluations also
identify the minimal simply connected nonspherical outputs and
weight-$d$ Engel manifolds, and the free OS49 Gysin lattice retains
the extra torsion outside the coprime sector
(\cref{maps:src:D.2.28}, \cref{maps:src:D.2.44}, \cref{maps:src:D.2.45}, \cref{maps:src:D.2.46}, \cref{maps:src:D.2.47}, \cref{maps:src:D.2.48}, \cref{night:P3F.34}).

\textit{Mechanism.} The primitive specialization and normal-circle-torus maps determine the integral cokernels before the full characteristic evaluations identify their smooth outputs.
\textit{Relation to literature.} The broader section-line construction and its realization question are credited to \cite[Remark~8.5]{E26}. The exact input table and the refined OS49 groups occur in the full statements below; the published classification theorems are cited there in the precise forms used.
\hyperlink{outline-III-F}{Proof in outline}; full statements: \SeriesRef{maps:src:D.2.28}, \SeriesRef{maps:src:D.2.44}, \SeriesRef{maps:src:D.2.45}, \SeriesRef{maps:src:D.2.46}, \SeriesRef{maps:src:D.2.47}, \SeriesRef{maps:src:D.2.48}, \SeriesRef{night:P3F.34}.

\hypertarget{main-III-G}{}
\paragraph{Theorem G (The conifold transition).}
\label{series:III:G}
Moving the zero of the line-bundle linearization gives the local
equation $xy=u(zv-s)$ and a proper analytic family with one ordinary
double point. The actual replacement of $S^2\times D^4$ by
$D^3\times S^3$ and its primitive boundary maps identify its two
smooth sides as the original six-sphere and $S^3\times S^3$,
the latter with algebraic dimension one
(\cref{maps:src:D.3.1}, \cref{maps:src:D.3.2}, \cref{maps:src:D.3.3}).

\textit{Mechanism.} The actual annular gluing gives the node equation and the primitive boundary maps of its surgery, determining the smooth side with its complex orientation.
\textit{Relation to literature.} The smooth classification and $\Theta_6=0$ are known inputs, cited in \cref{maps:src:D.3.2}; the full statements give the holomorphic construction and its marked surgery maps.
\hyperlink{outline-III-G}{Proof in outline}; full statements: \SeriesRef{maps:src:D.3.1}, \SeriesRef{maps:src:D.3.2}, \SeriesRef{maps:src:D.3.3}.

\subsection{The proofs in outline}
\hypertarget{outline-III-A}{}
\subsubsection*{Theorem A: The prescribed quotient and relative extensions}
\begin{enumerate}[label=\arabic*.,leftmargin=*]
\item\hypertarget{step-III-A-1}{} Use the primitive circle quotient and its explicit period
comparison from Paper~I.  The comparison preserves the signed primitive
covector, so the quotient and its boundary meridian are fixed maps.
\textit{Expansion:} \SeriesRef{mapsI:src:A.1.1}.
\item\hypertarget{step-III-A-2}{} The cusp target is the solid torus supplied by the local quotient
construction.  Its higher homotopy group in degree six is zero, which
gives the relative uniqueness within this prescribed target.
\textit{Expansion:} \SeriesRef{mapsI:src:A.2.1}.
\item\hypertarget{step-III-A-3}{} Choose a closed six-ball in the interior of the cusp
neighbourhood.  Collapse its boundary to obtain the relative pinch
$X\to X\vee S^6$, equal to the identity off the ball.
\textit{Expansion:} \SeriesRef{maps:src:A.2.2}.
\item\hypertarget{step-III-A-4}{} Compose the first summand with $e$ and the second with a
representative of $\alpha\in\pi_6(S^3)$.  The resulting map agrees
with $e$ on the exterior and represents $[e]+\alpha$.
This proves the exact scope of the extension assertion.
\textit{Expansion:} \SeriesRef{maps:src:A.2.2}.
\end{enumerate}

\hypertarget{outline-III-B}{}
\subsubsection*{Theorem B: The quotient and algebraic-reduction homotopy classes}
\begin{enumerate}[label=\arabic*.,leftmargin=*]
\item\hypertarget{step-III-B-1}{} Express the framing difference on the regular four-torus.
Its degree-one obstruction is the common invariant covector $\gamma$
modulo two.  Restriction to $\ker\gamma$ therefore lifts this
comparison to $\operatorname{Spin}$.
\textit{Expansion:} \SeriesRef{maps:src:B.1.1}.
\item\hypertarget{step-III-B-2}{} Twice the remaining degree-three obstruction is the restriction
of the common invariant class dual to the invariant vector $d$.
Since $d\in\ker\gamma$, this restriction is zero.  The integral
third cohomology of the three-torus is torsion-free, so the obstruction
itself is zero.
\textit{Expansion:} \SeriesRef{maps:src:B.1.1}.
\item\hypertarget{step-III-B-3}{} Pontryagin--Thom now represents $e$ by the Lie-framed
three-torus.  The periodic flat Dirac operator has paired nonzero
spectrum and a two-dimensional kernel.  The normalization in
\cite[Theorem~4.14, p.~423]{mapsAPS2} gives real $e$-invariant $1/2$.
\textit{Expansion:} \SeriesRef{maps:src:B.1.2}.
\item\hypertarget{step-III-B-4}{} Toda's stabilization calculation identifies the injected
$\pi_6(S^3)$ with the subgroup generated by $2\nu$ in the stable
three-stem.  The class with this invariant is $6\nu'$, also the
displayed composite of the three Hopf maps
\cite[Theorem~4.1 and Corollaries~4.2--4.3]{mapsToda52}.
\textit{Expansion:} \SeriesRef{maps:src:B.1.2}.
\item\hypertarget{step-III-B-5}{} In the two multiplicative coordinates, the normalized power map
is $P_{3,4}(z_1,z_2)=(z_1^3,z_2^4)/\|(z_1^3,z_2^4)\|$.
A regular value has twelve positively oriented inverse images, so its
degree is twelve.
\textit{Expansion:} \SeriesRef{maps:src:B.2.1}.
\item\hypertarget{step-III-B-6}{} The marked quotient construction gives the equality of maps
$f=H\circ P_{3,4}\circ e$, where $H:S^3\to S^2$ is the Hopf
map and $\PP^1$ is identified with $S^2$ by the stated orientation.
\textit{Expansion:} \SeriesRef{maps:src:B.2.1}.
\item\hypertarget{step-III-B-7}{} The degree-twelve self-map acts by multiplication by twelve on
$\pi_6(S^3)=\Z/12$.  Thus $P_{3,4}\circ e$ is null-homotopic;
postcomposition with $H$ gives the required null-homotopy of $f$.
The body also gives the normalized-section homotopy directly.
\textit{Expansion:} \SeriesRef{maps:src:B.2.1}.
\end{enumerate}

\hypertarget{outline-III-C}{}
\subsubsection*{Theorem C: The retained-angle cusp map and finite cell model}
\begin{enumerate}[label=\arabic*.,leftmargin=*]
\item\hypertarget{step-III-C-1}{} Define the map in the triple-point chart using the retained
angle and the two radial coordinates.  Keep the bounded period
correction when differentiating the normal frame.
\textit{Expansion:} \SeriesRef{maps:src:B.3.1}.
\item\hypertarget{step-III-C-2}{} The resulting frame differs from the product frame by
$\operatorname{diag}(e^{-i(a+b)},e^{ia},e^{ib})$ in
$\operatorname{SU}(3)$.  The displayed two $\operatorname{SU}(2)$
factor homotopies contract this comparison.
\textit{Expansion:} \SeriesRef{maps:src:B.3.1}.
\item\hypertarget{step-III-C-3}{} The inverse-image torus therefore has its Lie framing.  Its
spin quadratic form has Arf invariant one, so its Pontryagin--Thom
class is the nonzero $\eta_4\eta_5$.
\textit{Expansion:} \SeriesRef{maps:src:B.3.1}.
\item\hypertarget{step-III-C-4}{} Two allowed real-period lifts agree on the exterior collar;
their straight-line interpolation gives a relative homotopy.
The core-collapse comparison factors through the specified
$S^3\vee S^2$ cofiber, with the map using its $S^3$ summand.
These are the explicit comparisons among the cusp maps.
\textit{Expansion:} \SeriesRef{maps:src:B.3.2}, \SeriesRef{maps:src:B.3.3}.
\item\hypertarget{step-III-C-5}{} Construct the regular mapping-cylinder bundle using the actual
integral monodromies.  In every finite affine cap retain the drift
from the logarithmic translation; it is part of the attaching map. In the cusp write the seam map as
$[a,b,\theta]\mapsto[a,b-\theta(B_0a-v(a))]$.
The frozen-period comparison is an explicit logarithmic-coordinate
homeomorphism and agrees with the prescribed boundary markings.
\textit{Expansion:} \SeriesRef{maps:src:B.5.1}.
\item\hypertarget{step-III-C-6}{} The positive quotient has interior $T^2\times(0,1)$ and the
two marked boundary tori.  Its incompressible boundary and the
induced fundamental-group isomorphism meet Waldhausen's
boundary-fixed product criterion
\cite[Theorem~6.1, p.~77]{mapsWaldhausen68}. Triangulate these marked pieces and their common faces.  The
resulting mapping-cylinder cells have the displayed counts, and
their alternating sum is two.
\textit{Expansion:} \SeriesRef{maps:src:B.5.1}.
\item\hypertarget{step-III-C-7}{} Van Kampen gives the relations $x^3=x^4=1$.  The primitive
finite pullback lattices and the two cusp sweeps give unimodular
integral attachments, so $H_1=H_2=0$.  Duality, the Euler number,
Hurewicz and Whitehead finish the sphere identification.
\textit{Expansion:} \SeriesRef{maps:src:B.5.2}.
\end{enumerate}

\hypertarget{outline-III-D}{}
\subsubsection*{Theorem D: The relative Kummer cover and its global completion}
\begin{enumerate}[label=\arabic*.,leftmargin=*]
\item\hypertarget{step-III-D-1}{} On the $N$th root of the base use the actual period
$C(v^N)/N$.  The monomial maps and their lattice inclusions
separate the absolute degree-$N^3$ map from its relative
degree-$N^2$ map.
\textit{Expansion:} \SeriesRef{maps:src:C.1.1}.
\item\hypertarget{step-III-D-2}{} Across adjacent unimodular triangles, smoothness imposes the
displayed divisibility relations on primitive ray indices.
The reverse relations force equality, so a smooth full-fan cover
must use the scalar toric sublattice $n\Z^3$.
\textit{Expansion:} \SeriesRef{maps:src:C.1.2}.
\item\hypertarget{step-III-D-3}{} Continue the prescribed local Kummer lattice by the marked
order-four monodromy.  Its toric intersection is
$\Z w\oplus N\Z d$.  Its primitive ray indices violate that
scalar condition for $N>1$.
\textit{Expansion:} \SeriesRef{maps:src:C.1.3}.
\item\hypertarget{step-III-D-4}{} The explicit word
$A_2^2A_1A_2^2A_1^{-1}$ also produces $N^2$ distinct based
lattices; their residues are exhibited in the body.  The finite
semidirect quotient supplies a global cover with the required
local lattice.
\textit{Expansion:} \SeriesRef{maps:src:C.1.4}.
\item\hypertarget{step-III-D-5}{} Its normal monomial completion has a free basis indexed by
fractional exponents in a half-open fundamental box.  This proves
finite flatness of that normal completion while recording the
smoothness obstruction above.
\textit{Expansion:} \SeriesRef{maps:src:C.1.5}.
\end{enumerate}

\hypertarget{outline-III-E}{}
\subsubsection*{Theorem E: The cyclic universal covers}
\begin{enumerate}[label=\arabic*.,leftmargin=*]
\item\hypertarget{step-III-E-1}{} The connected cusp cover has $d$ components in its periodic
normalization.  Its Laurent relation matrix reduces integrally
to $\operatorname{diag}(1,t-1,t-1,0)$.
\textit{Expansion:} \SeriesRef{maps:src:D.1.1}.
\item\hypertarget{step-III-E-2}{} The nontrivial deck characters make the regular torus Koszul
complex acyclic.  The primitive finite-cap degree-two pullbacks
and the cusp sweep with coefficient $d$ determine the integral
groups, including their torsion.
\textit{Expansion:} \SeriesRef{maps:src:D.1.4}.
\item\hypertarget{step-III-E-3}{} The $d$ cusp divisors give classes $x_i$ with
$\sum x_i=0$ and
$x_ix_j=(\mathbf1_{j=i+1}+\mathbf1_{j=i-1}-2\mathbf1_{j=i})\kappa$,
where $\kappa$ generates the degree-four torsion.
\textit{Expansion:} \SeriesRef{maps:src:D.1.6}.
\item\hypertarget{step-III-E-4}{} For the finite-coefficient class $u$ the body computes
$\langle ux_ix_j,[M_d]\rangle$ by the same cyclic matrix,
$\langle u^2x_i,[M_d]\rangle=-12$, and
$\langle u^3,[M_d]\rangle=0$ modulo $d$.  The support-function
restriction has square $-12$, which fixes the second value.
\textit{Expansion:} \SeriesRef{maps:src:D.1.9}.
\item\hypertarget{step-III-E-5}{} Transfer kills the integral pullback in degree four, so
$c_2=p_1=0$.  Multiplication by four is injective on
$\Z^{d-1}\oplus\Z/d$, giving zero refined Pontryagin class.
The full $H_6(K(H_2,2))$ calculation shows that these evaluations
specify the remaining invariant; Zhubr's classification applies
to the already constructed covers.
\textit{Expansion:} \SeriesRef{maps:src:D.1.11}.
\item\hypertarget{step-III-E-6}{} Euler multiplicativity gives $\chi=2d$.  A connected sum of
$g$ copies of $S^3\times S^3$ has Euler number $2-2g$; for
$d>1$ the equality would require $g=1-d<0$.  This supplies the
numerical obstruction to that proposed connected-sum type.
\textit{Expansion:} \SeriesRef{maps:src:D.1.4}.
\end{enumerate}

\hypertarget{outline-III-F}{}
\subsubsection*{Theorem F: The marked elliptic profiles and their integral outputs}
\begin{enumerate}[label=\arabic*.,leftmargin=*]
\item\hypertarget{step-III-F-1}{} Fix the initial section pair, the component on which the
linearization vanishes, the marked filling and every affine
logarithmic vector.  The profile is the elliptic-fibre configuration;
a pattern is the entire marked initial row, with its stated
arithmetic conditions.
\textit{Expansion:} \SeriesRef{maps:src:D.2.28}.
\item\hypertarget{step-III-F-2}{} Compute the primitive specialization lattices and the actual
normal-circle-torus maps.  Regular-fibre restriction is carried
through the boundary and Gysin maps before the finite-coefficient
classes are used.  Their integral cokernels retain all torsion.
\textit{Expansion:} \SeriesRef{maps:src:D.2.28}.
\item\hypertarget{step-III-F-3}{} The all-order calculations for the original, III-star,
Liu--Rollenske and Wang families give their explicit determinant
or normal-torus matrices.  The additive and multiplicative local
calculations close the remaining profiles; the complete table in
\cref{maps:src:D.2.28} records every pattern.
\textit{Expansion:} \SeriesRef{maps:src:D.2.28}.
\item\hypertarget{step-III-F-4}{} In the simply connected residual cases retain the spin
characteristic class, Zhubr's refined class and the full
fundamental-class image.  The conic surgery and the marked
$Z_{c,\epsilon}$ surgeries give explicit representatives of the
types in \cref{maps:src:D.2.44}, \cref{maps:src:D.2.45}.
\textit{Expansion:} \SeriesRef{maps:src:D.2.44}, \SeriesRef{maps:src:D.2.45}.
\item\hypertarget{step-III-F-5}{} For a period sublattice of the cusp, intersect the component
restrictions on each degree-six del Pezzo normalization.  Keep
loops, repeated edges and repeated vertices in this calculation;
the complete trilinear expression is
\cref{maps:src:D.2.47}.
\textit{Expansion:} \SeriesRef{maps:src:D.2.47}.
\item\hypertarget{step-III-F-6}{} In the simply connected weight-$d$ sector the Seifert
determinant is $P=\pm1$.  The full invariant includes the
Pontryagin square modulo $d\gcd(d,2)$ and the third power
modulo $d\gcd(d,3)$, with evaluations $-12$ and $0$.
Together with the cyclic divisor matrix and refined class zero,
these give the single type $M_d$ in \cref{maps:src:D.2.48}.
\textit{Expansion:} \SeriesRef{maps:src:D.2.48}.
\item\hypertarget{step-III-F-7}{} For a free OS49 cap $(m,n,d)=(3k,n,d)$ with $n\ne0$,
put $g=\gcd(3k,n,3d)$.  The exact primitive Gysin lattice has
Smith invariants $g,|n|$, giving
$H_2=H_3=\Z/|n|\oplus\Z/g$ with $g\in\{1,3\}$.
For example $(9,3,1)$ gives $(\Z/3)^2$.
The zero-$n$ calculation and the precise freeness conditions are
retained in \cref{night:P3F.33}, \cref{night:P3F.34}.
\textit{Expansion:} \SeriesRef{night:P3F.33}, \SeriesRef{night:P3F.34}.
\end{enumerate}

\hypertarget{outline-III-G}{}
\subsubsection*{Theorem G: The conifold transition}
\begin{enumerate}[label=\arabic*.,leftmargin=*]
\item\hypertarget{step-III-G-1}{} Start with the marked elliptic surface and the line
$M=\OO(P-O)$.  Move the linearization as $\phi_s=\phi_0-s\phi_1$,
using its exact annular identification on the common exterior.
\textit{Expansion:} \SeriesRef{maps:src:D.3.1}.
\item\hypertarget{step-III-G-2}{} Near the collision the equation is
$xy=u(zv-s)$.  At $s=0$ this is the ordinary double point;
its small resolution has exceptional normal bundle
$\OO(-1)\oplus\OO(-1)$.  The root caps remain in the same
proper family because the collars were chosen uniformly.
\textit{Expansion:} \SeriesRef{maps:src:D.3.1}.
\item\hypertarget{step-III-G-3}{} On real manifolds the transition replaces $S^2\times D^4$
by $D^3\times S^3$.  General position makes the complement
simply connected; Alexander duality and the two primitive boundary
maps give $H_3=\Z^2$ and $H_2=H_4=0$.
\textit{Expansion:} \SeriesRef{maps:src:D.3.2}.
\item\hypertarget{step-III-G-4}{} The spin Wall invariant has the standard rank-two skew form
and zero normal invariant.  The classification, together with
$\Theta_6=0$, identifies the smoothing with $S^3\times S^3$.
\textit{Expansion:} \SeriesRef{maps:src:D.3.2}.
\item\hypertarget{step-III-G-5}{} The root canonical formula and the varying periods give
algebraic dimension one.  The formal cohomology of the exceptional
curve transports the stated structure-sheaf cohomology across the
small resolution.  The comparison with the published CDP quotient
uses its actual algebraic-reduction fibres in
\cref{maps:src:D.3.4}.
\textit{Expansion:} \SeriesRef{maps:src:D.3.3}, \SeriesRef{maps:src:D.3.4}.
\end{enumerate}

